\documentclass[10pt,aspectratio=169]{beamer}
\usepackage{bma01}

\title[Asíntotas verticales]{Límites infinitos y asíntotas verticales}
\subtitle{Ejemplos típicos: funciones racionales\\[2pt]\small Semana 4 -- Sesión 1}

\begin{document}

\begin{frame}
  \titlepage
\end{frame}

\begin{frame}{Objetivos de la sesión}
  \begin{enumerate}
    \item Definir $\lim_{x\to a}f(x)=\pm\infty$ y los límites laterales infinitos.
    \item Usar el \resaltar{análisis de signos} para determinar si $f\to+\infty$ ó $-\infty$.
    \item Hallar las \resaltar{asíntotas verticales} de funciones racionales, logarítmicas y trigonométricas.
    \item Distinguir una asíntota vertical de un ``hueco'' (discontinuidad evitable).
  \end{enumerate}
  \vspace{1mm}
  {\small\textbf{Competencia 2 del sílabo:} Evalúa propiedades de límites y continuidad y su utilidad en aplicaciones de ciencia e ingeniería.}
  \begin{block}{Plan (2 horas; incluye verificación previa de 5 min)}
    \begin{tabular}{ll}
      0:00 -- 0:25 & Límites infinitos: idea y definición\\
      0:25 -- 0:45 & Reglas para $\frac{L}{0}$; análisis de signos\\
      0:45 -- 1:30 & Asíntotas verticales: racionales, logaritmo, tangente\\
      1:30 -- 1:45 & Aplicaciones físicas\\
      1:45 -- 2:00 & Actividad colaborativa, resumen y cierre metacognitivo
    \end{tabular}
  \end{block}
\end{frame}

\begin{frame}{Aula invertida: preparación y verificación}
  \small
  \begin{previa}[antes de la sesión]
    \begin{enumerate}
      \item Lea límites infinitos y asíntotas verticales.
      \item Estudie el signo de $\frac{x+1}{(x-2)^2}$ cerca de $x=2$.
      \item Pregunta guía: ¿una función racional siempre tiene asíntota donde el denominador se anula?
    \end{enumerate}
  \end{previa}
  \begin{block}{Verificación al inicio (5 min, en pares)}
    \begin{enumerate}
      \item Calcule $\lim_{x\to0^+}\frac1x$ y $\lim_{x\to0^-}\frac1x$.
      \item ¿Qué diferencia hay entre asíntota y hueco?
    \end{enumerate}
    Los pares discuten y votan; el docente aclara solo lo que genera dudas y ajusta el ritmo de la sesión.
  \end{block}
\end{frame}

\section{Límites infinitos}

\begin{frame}{Idea: valores que crecen sin cota}
  \begin{columns}
    \begin{column}{0.5\textwidth}
      \begin{center}
        \small
        \begin{tabular}{c|c}
          $x$ & $\frac1{x^2}$\\ \hline
          $\pm0.1$ & $100$\\
          $\pm0.01$ & $10\,000$\\
          $\pm0.001$ & $1\,000\,000$
        \end{tabular}
      \end{center}
      Cuando $x\to0$, $\frac1{x^2}$ supera cualquier número: \ $\displaystyle\lim_{x\to0}\frac1{x^2}=+\infty$.
      \begin{definicion}
        $\displaystyle\lim_{x\to a}f(x)=+\infty$ si
        \[ \forall M>0\ \exists\delta>0:\ 0<|x-a|<\delta\Rightarrow f(x)>M. \]
        Análogo para $-\infty$ ($f(x)<-M$) y para límites laterales.
      \end{definicion}
    \end{column}
    \begin{column}{0.5\textwidth}
      \centering
      \begin{tikzpicture}
        \begin{axis}[bma, xmin=-2, xmax=2, ymin=-0.5, ymax=10, width=6.6cm, height=6cm, xtick={-1,1}, ytick={4}, yticklabels={$M$}, samples=200]
          \fill[uniVerde!20] (axis cs:-0.5,0) rectangle (axis cs:0.5,10);
          \addplot[curva, domain=-2:-0.31] {1/x^2};
          \addplot[curva, domain=0.31:2] {1/x^2};
          \draw[uniGranate, dashed] (axis cs:-2,4) -- (axis cs:2,4);
          \node at (axis cs:0,9) {\scriptsize $|x|<\delta$};
        \end{axis}
      \end{tikzpicture}

      \small Si $|x|<\delta=\frac{1}{\sqrt M}$, entonces $\frac1{x^2}>M$.
    \end{column}
  \end{columns}
  \small \alert{Nota:} ``$=+\infty$'' indica \emph{cómo} no existe el límite (como número real).
\end{frame}

\begin{frame}{Límites laterales infinitos}
  \begin{columns}
    \begin{column}{0.5\textwidth}
      \begin{ejemplo}[1]
        \[ \lim_{x\to0^+}\frac1x=+\infty,\qquad \lim_{x\to0^-}\frac1x=-\infty. \]
        No existe $\lim_{x\to0}\frac1x$ (ni siquiera como $\pm\infty$).
      \end{ejemplo}
      \begin{ejemplo}[2]
        $\displaystyle\lim_{x\to3^+}\frac{2}{x-3}=+\infty$ \ ($x-3\to0^+$), \quad $\displaystyle\lim_{x\to3^-}\frac{2}{x-3}=-\infty$.
      \end{ejemplo}
      \begin{ejemplo}[3]
        $\displaystyle\lim_{x\to0^+}\ln x=-\infty$.
      \end{ejemplo}
    \end{column}
    \begin{column}{0.5\textwidth}
      \centering
      \begin{tikzpicture}
        \begin{axis}[bma, xmin=-3, xmax=3, ymin=-5, ymax=5, width=6.6cm, height=6cm, xtick={-2,2}, ytick={-2,2}, samples=200]
          \addplot[curva, domain=-3:-0.2] {1/x};
          \addplot[curva, domain=0.2:3] {1/x};
          \node[uniAzul] at (axis cs:1.8,1.8) {\scriptsize $y=\frac1x$};
          \draw[->, uniGranate] (axis cs:0.5,3) -- (axis cs:0.25,4.6);
          \draw[->, uniGranate] (axis cs:-0.5,-3) -- (axis cs:-0.25,-4.6);
          \node[uniGranate] at (axis cs:1.2,4.3) {\scriptsize $+\infty$};
          \node[uniGranate] at (axis cs:-1.2,-4.3) {\scriptsize $-\infty$};
        \end{axis}
      \end{tikzpicture}
    \end{column}
  \end{columns}
\end{frame}

\begin{frame}{Reglas: ``número sobre cero''}
  \begin{teorema}
    Si $\lim_{x\to a}f(x)=L\neq0$ y $\lim_{x\to a}g(x)=0$, entonces $\left|\dfrac{f(x)}{g(x)}\right|\to+\infty$. El signo se determina así:
    \begin{center}
      \renewcommand{\arraystretch}{1.3}
      \begin{tabular}{c|c|c}
        $L$ & $g(x)\to0$ por & $\frac fg\to$\\ \hline
        $>0$ & $0^+$ & $+\infty$\\
        $>0$ & $0^-$ & $-\infty$\\
        $<0$ & $0^+$ & $-\infty$\\
        $<0$ & $0^-$ & $+\infty$
      \end{tabular}
    \end{center}
  \end{teorema}
  \begin{proposicion}[otras reglas]
    Si $f\to+\infty$ y $g\to L$ (finito): $f+g\to+\infty$; \ $fg\to+\infty$ si $L>0$ y $-\infty$ si $L<0$. \
    $(+\infty)+(+\infty)=+\infty$, \ $(+\infty)(+\infty)=+\infty$.
  \end{proposicion}
  \small \alert{Indeterminadas:} $\infty-\infty$, \ $0\cdot\infty$, \ $\frac\infty\infty$, \ $\frac00$. \ Requieren transformar la expresión.
\end{frame}

\begin{frame}{Análisis de signos: ejemplo}
  \begin{ejemplo}[4]
    $\displaystyle f(x)=\frac{x+1}{x^2-4}=\frac{x+1}{(x-2)(x+2)}$. Estudiemos $x\to2^\pm$ y $x\to-2^\pm$.
    \begin{center}
      \renewcommand{\arraystretch}{1.3}
      \begin{tabular}{c|c|c|c|c}
        límite & numerador & $(x-2)$ & $(x+2)$ & resultado\\ \hline
        $x\to2^+$ & $3$ & $0^+$ & $4$ & $+\infty$\\
        $x\to2^-$ & $3$ & $0^-$ & $4$ & $-\infty$\\
        $x\to-2^+$ & $-1$ & $-4$ & $0^+$ & $+\infty$\\
        $x\to-2^-$ & $-1$ & $-4$ & $0^-$ & $-\infty$
      \end{tabular}
    \end{center}
    Por ejemplo, en $x\to-2^+$: \ $\dfrac{-1}{(-4)(0^+)}=\dfrac{-1}{0^-}=+\infty$.
  \end{ejemplo}
  \small Truco: evalúe el signo con un número cercano, p.\,ej.\ $x=-1.99$: \ $\frac{-0.99}{(-3.99)(0.01)}>0$.
\end{frame}

\section{Asíntotas verticales}

\begin{frame}{Asíntota vertical}
  \begin{columns}
    \begin{column}{0.52\textwidth}
      \begin{definicion}
        La recta $x=a$ es una \resaltar{asíntota vertical} de la gráfica de $f$ si al menos uno de los límites laterales
        $\lim_{x\to a^\pm}f(x)$ es $+\infty$ ó $-\infty$.
      \end{definicion}
      \begin{block}{Dónde buscarlas}
        \begin{itemize}
          \item Ceros del denominador de una racional (que no se cancelen con el numerador).
          \item Bordes del dominio de $\ln$: $\ln g(x)$ con $g(x)\to0^+$.
          \item $\tan x$, $\sec x$ en $\frac\pi2+k\pi$; \ $\ctg x$, $\csc x$ en $k\pi$.
        \end{itemize}
      \end{block}
    \end{column}
    \begin{column}{0.48\textwidth}
      \centering
      \begin{tikzpicture}
        \begin{axis}[bma, xmin=-4.5, xmax=4.5, ymin=-5, ymax=5, width=6.8cm, height=6cm, xtick={-2,2}, ytick=\empty, samples=300, restrict y to domain=-8:8, unbounded coords=jump]
          \addplot[curva, domain=-4.5:4.5] {(x+1)/(x^2-4)};
          \draw[uniGranate, dashed] (axis cs:2,-5) -- (axis cs:2,5);
          \draw[uniGranate, dashed] (axis cs:-2,-5) -- (axis cs:-2,5);
          \node[uniGranate] at (axis cs:2.6,-4.3) {\scriptsize $x=2$};
          \node[uniGranate] at (axis cs:-2.8,4.3) {\scriptsize $x=-2$};
        \end{axis}
      \end{tikzpicture}

      \small $f(x)=\frac{x+1}{x^2-4}$ (Ejemplo 4).
    \end{column}
  \end{columns}
\end{frame}

\begin{frame}{Racionales: asíntota o hueco}
  \begin{columns}
    \begin{column}{0.55\textwidth}
      \begin{ejemplo}[5]
        $f(x)=\dfrac{x^2-3x+2}{x^2-4}=\dfrac{(x-1)(x-2)}{(x-2)(x+2)}$.
        \begin{itemize}
          \item En $x=2$ el factor se cancela: $\lim_{x\to2}f(x)=\lim_{x\to2}\frac{x-1}{x+2}=\frac14$. \ \textbf{Hueco} en $(2,\frac14)$, no asíntota.
          \item En $x=-2$: $\frac{-3}{0}$ \ $\Rightarrow$ \ asíntota vertical $x=-2$.
        \end{itemize}
      \end{ejemplo}
      \begin{teorema}
        Si $R=\frac PQ$ está \textbf{simplificada} (sin factores comunes), sus asíntotas verticales son exactamente las rectas $x=a$
        con $Q(a)=0$.
      \end{teorema}
    \end{column}
    \begin{column}{0.45\textwidth}
      \centering
      \begin{tikzpicture}
        \begin{axis}[bma, xmin=-5, xmax=5, ymin=-5, ymax=6, width=6.6cm, height=6cm, xtick={-2,2}, ytick={1}, samples=300, restrict y to domain=-8:8, unbounded coords=jump]
          \addplot[curva, domain=-5:5] {(x-1)/(x+2)};
          \draw[uniGranate, dashed] (axis cs:-2,-5) -- (axis cs:-2,6);
          \addplot[only marks, mark=o, uniAzul, fill=white, mark size=2.5pt] coordinates {(2,0.25)};
        \end{axis}
      \end{tikzpicture}
    \end{column}
  \end{columns}
\end{frame}

\begin{frame}{Más ejemplos}
  \begin{columns}
    \begin{column}{0.33\textwidth}
      \centering
      \begin{tikzpicture}
        \begin{axis}[bma, xmin=-2.5, xmax=2.5, ymin=-0.5, ymax=8, width=5cm, height=4.6cm, xtick={-1,1}, ytick=\empty, samples=300]
          \addplot[curva, domain=-2.5:-1.06] {1/(x^2-1)^2};
          \addplot[curva, domain=-0.94:0.94] {1/(x^2-1)^2};
          \addplot[curva, domain=1.06:2.5] {1/(x^2-1)^2};
          \draw[uniGranate, dashed] (axis cs:-1,0) -- (axis cs:-1,8);
          \draw[uniGranate, dashed] (axis cs:1,0) -- (axis cs:1,8);
        \end{axis}
      \end{tikzpicture}\\
      \small $\frac{1}{(x^2-1)^2}$: $+\infty$ por ambos lados (factor al cuadrado).
    \end{column}
    \begin{column}{0.33\textwidth}
      \centering
      \begin{tikzpicture}
        \begin{axis}[bma, xmin=-0.5, xmax=4.5, ymin=-3.5, ymax=2.5, width=5cm, height=4.6cm, xtick={1,3}, ytick=\empty, samples=200]
          \addplot[curva2, domain=1.01:2.99] {ln((x-1)*(3-x))};
          \draw[uniGranate, dashed] (axis cs:1,-3.5) -- (axis cs:1,2.5);
          \draw[uniGranate, dashed] (axis cs:3,-3.5) -- (axis cs:3,2.5);
        \end{axis}
      \end{tikzpicture}\\
      \small $\ln\big((x-1)(3-x)\big)$: dominio $(1,3)$; asíntotas $x=1$, $x=3$.
    \end{column}
    \begin{column}{0.33\textwidth}
      \centering
      \begin{tikzpicture}
        \begin{axis}[bma, xmin=-4.8, xmax=4.8, ymin=-4, ymax=4, width=5cm, height=4.6cm, xtick={-1.5708,1.5708}, xticklabels={$-\frac\pi2$,$\frac\pi2$}, ytick=\empty, samples=300, restrict y to domain=-6:6, unbounded coords=jump]
          \addplot[uniVerde, very thick, domain=-4.8:4.8] {tan(deg(x))};
          \draw[uniGranate, dashed] (axis cs:-4.7124,-4) -- (axis cs:-4.7124,4);
          \draw[uniGranate, dashed] (axis cs:-1.5708,-4) -- (axis cs:-1.5708,4);
          \draw[uniGranate, dashed] (axis cs:1.5708,-4) -- (axis cs:1.5708,4);
          \draw[uniGranate, dashed] (axis cs:4.7124,-4) -- (axis cs:4.7124,4);
        \end{axis}
      \end{tikzpicture}\\
      \small $\tan x$: infinitas asíntotas $x=\frac\pi2+k\pi$.
    \end{column}
  \end{columns}
  \begin{ejemplo}[6]
    $\displaystyle\lim_{x\to0^+}e^{1/x}=+\infty$, pero $\displaystyle\lim_{x\to0^-}e^{1/x}=0$: \ $x=0$ es asíntota vertical (basta \emph{un} lado).
  \end{ejemplo}
\end{frame}

\begin{frame}{Aplicaciones físicas}
  \begin{columns}
    \begin{column}{0.55\textwidth}
      \begin{ejemplo}[7: relatividad]
        El factor de Lorentz $\gamma(v)=\dfrac{1}{\sqrt{1-v^2/c^2}}$ cumple
        \[ \lim_{v\to c^-}\gamma(v)=+\infty. \]
        La energía $E=\gamma mc^2$ crece sin cota: ninguna partícula con masa alcanza la velocidad de la luz.
      \end{ejemplo}
      \begin{ejemplo}[8: lente delgada]
        $\frac1{d_o}+\frac1{d_i}=\frac1f\Rightarrow d_i=\dfrac{fd_o}{d_o-f}$. \ Si el objeto se acerca al foco ($d_o\to f^+$), la imagen se va al infinito.
      \end{ejemplo}
    \end{column}
    \begin{column}{0.45\textwidth}
      \centering
      \begin{tikzpicture}
        \begin{axis}[bma, xlabel={$v/c$}, ylabel={$\gamma$}, xmin=0, xmax=1.1, ymin=0, ymax=8, width=6.3cm, height=5.5cm, xtick={0.5,1}, ytick={1}, samples=200]
          \addplot[curva, domain=0:0.991] {1/sqrt(1-x^2)};
          \draw[uniGranate, dashed] (axis cs:1,0) -- (axis cs:1,8);
        \end{axis}
      \end{tikzpicture}
    \end{column}
  \end{columns}
\end{frame}

\begin{frame}{Actividad colaborativa: Diseñe la función}
  \small
  \begin{activa}[Aprendizaje inverso (de propiedades a fórmula)]
    Cada grupo diseña una función racional con asíntotas verticales en $x=-1$ y $x=3$, un \textbf{hueco} en $x=2$ y $f(x)>0$ para $x>3$. Intercambia su fórmula con otro grupo, que verifica cada propiedad con signos y con GeoGebra, y devuelve observaciones.
  \end{activa}
  \begin{itemize}
    \item \textbf{Organización:} equipos de 3 a 4 con roles rotativos (coordinador, redactor, verificador, portavoz); 12--15 minutos.
    \item \textbf{Producto:} Fórmula, tabla de signos y verificación gráfica del otro grupo.
    \item \textbf{Retroalimentación:} el docente recorre los equipos, hace preguntas y no da la solución hasta la puesta en común.
  \end{itemize}
\end{frame}

\section{Ejercicios}

\begin{frame}{Ejercicios propuestos}
  \small
  \begin{ejercicio}
    \begin{enumerate}
      \item Calcule: (a) $\displaystyle\lim_{x\to1^-}\frac{x+2}{x^2-1}$, \ (b) $\displaystyle\lim_{x\to3^+}\frac{x^2}{9-x^2}$, \ (c) $\displaystyle\lim_{x\to0^-}\left(\frac1x-\frac1{x^2}\right)$,
        \ (d) $\displaystyle\lim_{x\to\pi^-}\ctg x$, \ (e) $\displaystyle\lim_{x\to2^+}\frac{\ent{x}-x}{x-2}$.
      \item Halle todas las asíntotas verticales y el comportamiento a cada lado: (a) $f(x)=\dfrac{x^2+1}{x^2-x-6}$,
        \ (b) $g(x)=\dfrac{x^3-8}{x^2-4}$, \ (c) $h(x)=\dfrac{\sen x}{x(x-\pi)}$, \ (d) $k(x)=\ln|x^2-1|$.
      \item Construya una función racional con asíntotas verticales $x=-1$ y $x=3$, un hueco en $x=1$, y que tienda a $+\infty$ por ambos lados de $x=3$.
      \item Pruebe con la definición que $\displaystyle\lim_{x\to1}\frac{-2}{(x-1)^2}=-\infty$.
      \item ¿Es $x=0$ asíntota vertical de $f(x)=x\ln|x|$? ¿Y de $g(x)=\frac{1}{x}\sen x$?
      \item La ley de Coulomb $F=\frac{kq_1q_2}{r^2}$: describa el comportamiento cuando $r\to0^+$ para cargas de igual y de distinto signo.
    \end{enumerate}
  \end{ejercicio}
\end{frame}

\begin{frame}{Metacognición y uso responsable de la IA}
  \small
  \begin{columns}[t]
    \begin{column}{0.5\textwidth}
      \begin{metacog}[cierre, 3 min]
        \begin{enumerate}
          \metacomun
          \item Al diseñar en lugar de analizar, ¿qué comprendí mejor sobre el papel del denominador?
        \end{enumerate}
      \end{metacog}
    \end{column}
    \begin{column}{0.5\textwidth}
      \begin{iabox}
        Pida a una IA una función con las condiciones anteriores. Compruebe cada condición por su cuenta; si alguna falla, explique por qué.\\[2pt]
        \textit{Regla:} primero intente usted; luego use la herramienta para contrastar, nunca para sustituir su razonamiento. Anote el cambio en su bitácora.
      \end{iabox}
    \end{column}
  \end{columns}
\end{frame}

\begin{frame}{Resumen}
  \begin{itemize}
    \item $\lim_{x\to a}f=+\infty$: $f$ supera cualquier $M$ cerca de $a$.
    \item ``$\frac{L}{0}$'' con $L\ne0$: infinito; el signo por análisis de signos ($0^+$, $0^-$).
    \item $x=a$ es asíntota vertical si algún lateral es $\pm\infty$.
    \item Racionales simplificadas: asíntotas en los ceros del denominador; factores cancelados dan huecos.
    \item Otras fuentes: $\ln$ cerca de $0^+$, $\tan$, $\sec$, $e^{1/x}$.
  \end{itemize}
  \vspace{3mm}
  \textbf{Próxima sesión:} límites al infinito, asíntotas horizontales y oblicuas.
\end{frame}


\begin{frame}{Evaluación y trabajo posterior}
  \small
  \begin{evalbox}[Práctica del miércoles de la semana 4: PC1]
    \textbf{Cobertura:} semanas 1, 2 y 3. \textbf{Naturaleza:} práctica calificada individual escrita.
  \end{evalbox}
  \begin{itemize}
    \item \criterios
    \item \declaracionia
    \item \textbf{Portafolio:} guarde el producto de la actividad, la bitácora de metacognición y los ejercicios propuestos de esta sesión; son evidencia de progreso y base de la retroalimentación.
  \end{itemize}
\end{frame}

\end{document}
